\section{Frozen V4 confirmation attempt}
\label{app:v4-confirmation}

% CLAIM: C14,C15,C16,C17,C18,C19,E02
\subsection{Protocol, freeze, and evidence chain}

V4 was designed as a fresh, exactly-once, same-follower, whole-trajectory
confirmation attempt.  The frozen source, lock, manifest, protocol status, and
artifact index are registered separately from the V3 evidence chain in
\cref{app:provenance}.  The test identities and seeds had zero overlap
with V1--V3.  Test visibility occurred only after the lock and confirmation
authorization were committed.  The experiment then consumed the locked test
once, completed the primary raw bundle and statistical estimation, and froze
with status \texttt{\VFourProtocolStatus}.  No raw-stage resume or same-test
rerun is permitted.

\begin{table}[H]
  \centering
  \caption{V4 exactly-once timeline, commits, and hashes.  The report-only
  handoff follows the complete raw confirmation and does not authorize an
  experiment rerun.}
  \label{tab:v4-evidence-chain}
  \scriptsize
  % Generated; do not edit.
\begin{tabular}{@{}>{\raggedright\arraybackslash}p{0.14\linewidth}>{\raggedright\arraybackslash}p{0.34\linewidth}>{\raggedright\arraybackslash}p{0.42\linewidth}@{}}
\toprule
Stage & Frozen action/state & Commit or SHA-256 \\
\midrule
Protocol lock & exactly-once V4 protocol; base commit \texttt{1d5cba1b3e80\allowbreak{}72bcf2a9a404\allowbreak{}92e044d2af4c\allowbreak{}f9fe} & \texttt{EXPERIMENT\_PROTOCOL\_V4.md}: \texttt{baad38320593\allowbreak{}695a4c231f18\allowbreak{}02faa3a48b4a\allowbreak{}32b318da841f\allowbreak{}da5b1354cd8b\allowbreak{}770e} \\
Fresh split & locked V4 identities/seeds and zero-overlap split & \texttt{split\_manifest\_v4.json}: \texttt{1727505734c8\allowbreak{}026ed18d8712\allowbreak{}3d5d5a8c02e2\allowbreak{}f201a33ea786\allowbreak{}fbcde2c9ab39\allowbreak{}8796} \\
Configuration lock & frozen V4 method and execution configuration & \texttt{config\_lock\_v4.json}: \texttt{d61b0f8596b0\allowbreak{}4358c7bef6a1\allowbreak{}e43b6775b3db\allowbreak{}b00020c2aca2\allowbreak{}8d5d2cd4d9f6\allowbreak{}f3d3} \\
Pre-test status & \texttt{locked\_test\_unseen}; test identities unseen; same-test rerun prohibited & \texttt{protocol\_status\_v4.json}: \texttt{c0c3d358c969\allowbreak{}dbb343ac05dc\allowbreak{}964075a514f3\allowbreak{}7d8153ce47d6\allowbreak{}e4ca60a252de\allowbreak{}4909} \\
Raw confirmation & \texttt{confirm}; one locked invocation; execution count 1 & \texttt{461fc560461b\allowbreak{}0a4726cbabdb\allowbreak{}97b2dbd4dc30\allowbreak{}5e0a} \\
Post-test freeze & \VFourProtocolStatus; raw resume and same-test rerun prohibited & status artifact: \texttt{48c98a81a761\allowbreak{}29a0fc2dd913\allowbreak{}aabb28bc9312\allowbreak{}d31a76a4283b\allowbreak{}27bf1fea9431\allowbreak{}a34b} \\
Reporting repair & verified reporting-provenance repair; no raw experiment rerun & reporting: \texttt{8baece6b7051\allowbreak{}ccc231d9bb03\allowbreak{}62fd85e4aa5a\allowbreak{}94e5} \\
Bounded result & bounded V4 result and evidence package published & result: \texttt{f49b4ef1cacf\allowbreak{}8228c5d24335\allowbreak{}3184acb8a7d0\allowbreak{}2311} \\
Diagnostic aid & report-only same-information diagnosis; frozen gate unchanged & aid: \texttt{b9301eaf36dc\allowbreak{}04f1abf662c4\allowbreak{}2821eddfe8c3\allowbreak{}188a} \\
Root of trust & 152 bounded artifacts; all indexed hashes verified & \texttt{artifact\_index.json}: \texttt{fd78eb559d03\allowbreak{}9620ae1c6e06\allowbreak{}faac44ab6fc8\allowbreak{}dbff9208c055\allowbreak{}23b4efcab4a7\allowbreak{}5a95} \\
\bottomrule
\end{tabular}

\end{table}

The report-only repair did not call the generator or pipeline, change raw
artifacts, or recompute algorithm outputs.  V3 frozen bytes remained
unchanged, and the V4 evidence tree is likewise treated as immutable input to
the paper.

\subsection{Primary observation and complete gate disposition}

The primary analysis contains
\VFourPrimaryPairedTrajectoryCount/\VFourTestTrajectoryCount{} paired
trajectories and reports a \VFourPrimaryRelativeImprovement\% observed
relative RMSE reduction with 95\% CI
[\VFourPrimaryRelativeCILow\%, \VFourPrimaryRelativeCIHigh\%].  The observed
absolute reduction is \SI{\VFourPrimaryAbsoluteImprovement}{\radian}, with
95\% CI
[\VFourPrimaryAbsoluteCILow, \VFourPrimaryAbsoluteCIHigh]\,rad.  The
statistical label is \texttt{\VFourStatisticalClassification}, but the
effective label is \texttt{\VFourEffectiveClassification}.  The latter
controls the claim.  Because the preregistered validity gate failed, this
observed effect is retained as non-confirmatory rather than as a comparative
performance benefit.

\begin{table}[H]
  \centering
  \caption{Complete high-level disposition of the frozen V4 result.  The
  observed effect and all failed gates are shown together; the result is
  non-confirmatory.}
  \label{tab:v4-appendix-gates}
  \scriptsize
  % Generated; do not edit.
\begin{tabular}{@{}>{\raggedright\arraybackslash}p{0.22\linewidth}>{\raggedright\arraybackslash}p{0.45\linewidth}>{\raggedright\arraybackslash}p{0.17\linewidth}@{}}
\toprule
Gate or classification & Frozen observation & Status \\
\midrule
Fresh/exactly-once execution & one locked confirmation invocation followed by immutable report-only handoff & retained \\
Primary cycle/sample denominator & \num{\VFourPrimarySampleCountPerMethod} aligned cycles per direct method & pass \\
Paired trajectory denominator & \VFourPrimaryPairedTrajectoryCount/\VFourTestTrajectoryCount{} trajectories & pass \\
Observed primary effect & \VFourPrimaryRelativeImprovement\% relative RMSE reduction (95\% CI [\VFourPrimaryRelativeCILow\%, \VFourPrimaryRelativeCIHigh\%]) & \VFourStatisticalClassification \\
Direct method purity & \VFourDirectMethodPurityRate{} for P, PV, and PVA & pass \\
Same information & \VFourSameInformationFailureCount/\num{\VFourSameInformationAuditCycleCount} composite event-flag differences & fail \\
Safety/completion & \VFourPrimaryFailureCount{} failures; \VFourFallbackEventCount{} fallbacks; \VFourContinuousConstraintViolationCount{} violations & pass \\
Maximum-error guardrail & \VFourMaxErrorRelativeImprovement\% observed relative reduction & pass \\
Lag noninferiority & P/PVA mean lag \VFourPMeanLagMS/\VFourPVAMeanLagMS\,ms & fail: not established \\
Hard runtime & full instrumented Python pipeline; every direct method missed at least one frozen threshold & fail \\
Effective result & \VFourEffectiveClassification & non-confirmatory \\
\bottomrule
\end{tabular}

\end{table}

\subsection{Same-information diagnosis}

The shared measurement, estimator posterior, predictor output, horizon,
governor, follower, plant, initial-state policy, current-state policy, control
time, limits, and raw target position were audited over aligned cycles.
\VFourSameInformationFailureCount/
\num{\VFourSameInformationAuditCycleCount}{} composite event-flag entries
failed.  In every failed row, \texttt{deadline\_miss} was the only differing
token; all non-event shared fields and configuration-identity checks passed.
Thus the shared estimator and predictor fields passed.  Because the composite
audit was preregistered and frozen, its failure remains dispositive even
though the diagnosis is narrow.

\begin{table}[H]
  \centering
  \caption{All failed V4 same-information rows, generated from the bounded
  reviewer aid.  Each row differs only in the composite execution-time
  \texttt{deadline\_miss} token; the frozen gate remains failed.}
  \label{tab:v4-same-information-failures}
  \scriptsize
  % Generated; do not edit.
\begin{tabular}{@{}>{\raggedright\arraybackslash}p{0.31\linewidth}r>{\raggedright\arraybackslash}p{0.07\linewidth}>{\raggedright\arraybackslash}p{0.07\linewidth}>{\raggedright\arraybackslash}p{0.07\linewidth}>{\raggedright\arraybackslash}p{0.18\linewidth}@{}}
\toprule
Trajectory & $k$ & P DM & PV DM & PVA DM & Differing token \\
\midrule
\texttt{piecewise\_constant\_jerk/001} & 154 & -- & -- & yes & \texttt{deadline\_miss} \\
\texttt{piecewise\_constant\_jerk/003} & 86 & yes & -- & -- & \texttt{deadline\_miss} \\
\texttt{rapid\_reversal/015} & 215 & -- & yes & -- & \texttt{deadline\_miss} \\
\texttt{stationary\_endpoint/013} & 254 & yes & -- & -- & \texttt{deadline\_miss} \\
\texttt{stationary\_endpoint/014} & 260 & -- & yes & -- & \texttt{deadline\_miss} \\
\bottomrule
\end{tabular}

\end{table}

\subsection{Adverse cases, families, and contextual analyses}

PVA was harmful relative to P on
\VFourHarmfulCount/\VFourHarmfulDenominator{} primary trajectories.  Rapid
reversal was the weakest family, with observed relative estimate
\VFourRapidReversalRelativeImprovement\% and 95\% CI
[\VFourRapidReversalRelativeCILow\%,
\VFourRapidReversalRelativeCIHigh\%].  These trajectories and family-level
heterogeneity are retained and cannot be removed to improve the classification.

\begin{table}[H]
  \centering
  \caption{All frozen descriptive V4 family effects and the mechanically
  selected worst five trajectories.  Harmful rows and rapid-reversal
  heterogeneity remain part of the evidence; subgroups do not override the
  aggregate validity gate.}
  \label{tab:v4-family-effects}
  \scriptsize
  % Generated; do not edit.
\begin{tabular}{@{}>{\raggedright\arraybackslash}p{0.12\linewidth}>{\raggedright\arraybackslash}p{0.28\linewidth}>{\raggedright\arraybackslash}p{0.08\linewidth}>{\raggedright\arraybackslash}p{0.26\linewidth}>{\raggedright\arraybackslash}p{0.08\linewidth}@{}}
\toprule
Evidence & Family/trajectory & $n$/rank & Observed effect [95\% CI] & Harmful \\
\midrule
Family & Stationary endpoint & 20 & 92.96\% [91.53, 94.02] & 0/20 \\
Family & Oscillatory & 20 & 72.80\% [64.37, 78.46] & 1/20 \\
Family & Piecewise-constant jerk & 20 & 89.92\% [87.57, 91.87] & 0/20 \\
Family & Stop and go & 20 & 95.37\% [94.96, 95.74] & 0/20 \\
Family & Rapid reversal & 20 & 32.52\% [-15.73, 65.92] & 4/20 \\
Family & Boundary grazing & 20 & 36.64\% [32.46, 40.33] & 0/20 \\
Worst & \texttt{rapid\_reversal/007} & \#1 & $\Delta$RMSE = +0.0751\,rad & yes \\
Worst & \texttt{rapid\_reversal/010} & \#2 & $\Delta$RMSE = +0.0593\,rad & yes \\
Worst & \texttt{rapid\_reversal/005} & \#3 & $\Delta$RMSE = +0.0407\,rad & yes \\
Worst & \texttt{oscillatory/019} & \#4 & $\Delta$RMSE = +0.0261\,rad & yes \\
Worst & \texttt{rapid\_reversal/003} & \#5 & $\Delta$RMSE = +0.0074\,rad & yes \\
\bottomrule
\end{tabular}

\end{table}

The ordinary-\Ruckig{} comparison is contextual secondary evidence.  Its S5
predicted-P versus raw-predicted-PVA comparison has an incomplete PVA
denominator, so complete-pair inference is unavailable; neither complete-case
deletion nor substitution into the primary analysis is permitted.  The oracle
target-component matrix is an offline, noncausal, nondeployable diagnostic and
does not support an online or causal performance claim.

\begin{table}[H]
  \centering
  \caption{Frozen runtime measurements and contextual ordinary-\Ruckig{} and
  oracle evidence.  Runtime applies to the instrumented full Python pipeline;
  incomplete ordinary comparisons and noncausal oracle diagnostics do not
  replace the primary inference.}
  \label{tab:v4-runtime-context}
  \scriptsize
  % Generated; do not edit.
\begin{tabular}{@{}>{\raggedright\arraybackslash}p{0.22\linewidth}>{\raggedright\arraybackslash}p{0.12\linewidth}>{\raggedright\arraybackslash}p{0.26\linewidth}>{\raggedright\arraybackslash}p{0.22\linewidth}@{}}
\toprule
Evidence & Denominator & Observed & Status/boundary \\
\midrule
Runtime P & \num{\VFourRuntimeCycleCountPerMethod} & p99 \VFourPRuntimePNinetyNineUS\,$\mu$s; max \VFourPRuntimeMaxUS\,$\mu$s; \VFourPDeadlineMissCount{} misses & fail \\
Runtime PV & \num{\VFourRuntimeCycleCountPerMethod} & p99 \VFourPVRuntimePNinetyNineUS\,$\mu$s; max \VFourPVRuntimeMaxUS\,$\mu$s; \VFourPVDeadlineMissCount{} misses & fail \\
Runtime PVA & \num{\VFourRuntimeCycleCountPerMethod} & p99 \VFourPVARuntimePNinetyNineUS\,$\mu$s; max \VFourPVARuntimeMaxUS\,$\mu$s; \VFourPVADeadlineMissCount{} misses & fail \\
Ordinary Ruckig, raw PV & 116/120 & incomplete completion & contextual only; no paired inference \\
Ordinary Ruckig, raw PVA & 108/120 & incomplete completion & contextual only; no paired inference \\
Oracle PV vs. P & 120/120 & 1.40\% observed & offline noncausal diagnostic; nondeployable \\
Oracle PVA vs. PV & 120/120 & 79.38\% observed & offline noncausal diagnostic; nondeployable \\
Oracle PVA vs. PV, acceleration-active & 40/40 & 76.71\% observed & offline noncausal diagnostic; nondeployable \\
\bottomrule
\end{tabular}

\end{table}

\subsection{Safety, lag, runtime, and final boundary}

All three primary methods retained native direct execution and method purity
\VFourDirectMethodPurityRate{}.  The frozen synthetic protocol recorded
\VFourPrimaryFailureCount{} primary failures,
\VFourFallbackEventCount{} fallbacks, and
\VFourContinuousConstraintViolationCount{} continuous-constraint or
registered-invariant violations.  This is command-model evidence only, not
hardware or production-safety certification.

Max-error noninferiority passed, but lag noninferiority was not established;
the mean P and PVA lag point estimates were \VFourPMeanLagMS{} and
\VFourPVAMeanLagMS{} ms.  The instrumented full-Python pipeline failed the
hard-runtime gate.  P/PV/PVA p99 measurements were
\VFourPRuntimePNinetyNineUS/\VFourPVRuntimePNinetyNineUS/%
\VFourPVARuntimePNinetyNineUS\,\si{\micro\second}, maxima were
\VFourPRuntimeMaxUS/\VFourPVRuntimeMaxUS/%
\VFourPVARuntimeMaxUS\,\si{\micro\second}, and deadline-miss counts were
\VFourPDeadlineMissCount/\VFourPVDeadlineMissCount/%
\VFourPVADeadlineMissCount.  These measurements neither establish real-time
deployment readiness nor prove that an isolated or compiled implementation is
incapable of meeting the control period.

The frozen V4 data therefore contain a strong observed effect without
confirmatory validity.  They may be used descriptively only with the failed
gate and non-confirmatory boundary disclosed.  They do not establish PVA
superiority, may not be repaired through a same-test rerun, and provide no
real-stream, hardware, causal-oracle, or deployment confirmation.
